%% Chapter 1: boxes
\chapter{Boxes}
%
% ============================================================
% Numbered boxes
%   \begin{defbox}{Title}{label} ... \end{defbox}
%   - numbered per section, sharing one counter across all box types
%   - referenced with \ref{<prefix>:label}, e.g. \ref{def:vector_space}
%   - listed in the index at the end of the notes
%   - the starred form \begin{defbox*}{Title} is unnumbered and not indexed
%
%   Box         Title          Prefix   Needs
%   defbox      Definition     def
%   stbox       Setup          st
%   exbox       Example        ex
%   agrbox      Algorithm      agr
%   tecbox      Technique      tec
%   asmpbox     Assumption     asmp     (own counter, not in the index)
%   lembox      Lemma          lem      math
%   thmbox      Theorem        thm      math
%   prpbox      Proposition    prp      math
%   fctbox      Fact           fct      math
%   solbox      Optimization Solution  sol  math
%   codebox     Code Snippet   code     code
%
% Inside any box, \tcbline draws a separator (e.g. between a statement and its proof).
% \deftxt{term} typesets a newly defined term in bold italics.
% ============================================================
\section{Definitions and Examples}
%
\begin{defbox}{Vector Space}{vector_space}
A \deftxt{vector space} over $\R$ is a set $V$ with an addition $V\times V\to V$ and a scalar multiplication $\R\times V\to V$ such that:
\begin{itemize}
    \item $(V,+)$ is an abelian group;
    \item scalar multiplication is compatible with the field operations:
    \begin{itemize}
        \item $a(b\bs v)=(ab)\bs v$ and $1\bs v=\bs v$;
        \item $a(\bs u+\bs v)=a\bs u+a\bs v$ and $(a+b)\bs v=a\bs v+b\bs v$.
    \end{itemize}
\end{itemize}
\end{defbox}
%
\begin{exbox}{Euclidean Space}{euclid}
$\R^n$ with component-wise addition and scalar multiplication is a vector space in the sense of Definition \ref{def:vector_space}.
\end{exbox}
%
\begin{stbox}{Consumer Problem}{consumer}
A consumer with income $m>0$ faces prices $\bs p\in\R^n_{++}$ and solves
\[
\max_{\bs x\in\R^n_+} u(\bs x)\st \bs p\cdot\bs x\le m.
\]
\end{stbox}
%
\begin{defbox*}{An Unnumbered Definition}
The starred form has no number, no label and no index entry.
\end{defbox*}
%
\section{Results and Proofs}
%
\begin{lembox}{Triangle Inequality}{triangle}
For all $x,y\in\R$, $\abs{x+y}\le\abs{x}+\abs{y}$.
\end{lembox}
%
\begin{thmbox}{Reverse Triangle Inequality}{rev_triangle}
For all $x,y\in\R$,
\[
\abs[\big]{\abs{x}-\abs{y}}\le\abs{x-y}.
\]
\tcbline
\begin{proof}
By Lemma \ref{lem:triangle}, $\abs{x}=\abs{(x-y)+y}\le\abs{x-y}+\abs{y}$, so $\abs{x}-\abs{y}\le\abs{x-y}$. Swapping $x$ and $y$ gives $\abs{y}-\abs{x}\le\abs{x-y}$.
\end{proof}
\end{thmbox}
%
\begin{prpbox}{Solution of the Consumer Problem}{consumer_sol}
In Setup \ref{st:consumer}, if $u$ is continuous, a solution exists.
\tcbline
\begin{proof}
The budget set $\cb{\bs x\in\R^n_+:\bs p\cdot\bs x\le m}$ is closed and bounded, hence compact, and a continuous function attains its maximum on a compact set.
\end{proof}
\end{prpbox}
%
\begin{fctbox}{Geometric Series}{geom}
For $\abs{r}<1$, $\sum_{t=0}^\infty r^t=\frac{1}{1-r}$.
\end{fctbox}
%
\section{Other Boxes}
%
% asmpbox: numbered within each section with its own counter
\begin{asmpbox}{Preferences}{pref}
$u$ is continuous, strictly increasing and strictly quasi-concave.
\end{asmpbox}
%
\begin{agrbox}{Bisection}{bisection}
\begin{enumerate}
    \item Start with $a<b$ such that $f(a)f(b)<0$.
    \item Let $c=(a+b)/2$. If $f(a)f(c)<0$ set $b=c$, otherwise set $a=c$.
    \item Repeat until $b-a$ is below the tolerance.
\end{enumerate}
\end{agrbox}
%
\begin{tecbox}{Guess and Verify}{guess}
Guess a functional form with unknown coefficients, plug it into the equation and solve for the coefficients.
\end{tecbox}
%
% genbox: a general box with a free title, not numbered or indexed
\begin{genbox}{A General Box}
Use \texttt{genbox} for anything that does not fit the numbered types.
\end{genbox}
%
\begin{notebox}
Use \texttt{notebox} for remarks and warnings.
\end{notebox}
%
\begin{refbox}
Use \texttt{refbox} for quotes or references to other material.
\end{refbox}
